\section*{Bab \ref{ch:g12:quantum-world} --- Dunia Kuantum: Foton dan Tingkat Energi}

\begin{solution}{exo:g12:quantum-world:1}
$E = hc/\lambda = \num{1.99e-25}/\num{5.30e-7} = \qty{3.75e-19}{J} =
\qty{2.35}{eV}$. $N = \num{1.0e-3}/\num{3.75e-19} \approx \num{2.7e15}$
\omterm{def:g12:quantum-world:photon}{foton} tiap sekon.
\end{solution}

\begin{solution}{exo:g12:quantum-world:2}
$E = h\nu = 6.63\times10^{-34} \times 10^{8} = \qty{6.63e-26}{J}
\approx \qty{4.1e-7}{eV}$ --- sekitar lima juta kali lebih kecil daripada
\omterm{def:g12:quantum-world:photon}{foton} tampak. \omterm{def:g10:signals-and-waves:signal}{Sinyal} radio mana pun yang terdeteksi melibatkan begitu banyak \omterm{def:g12:quantum-world:photon}{foton}
tiap sekon sehingga alirannya mulus sempurna.
\end{solution}

\begin{solution}{exo:g12:quantum-world:3}
$\nu_0 = W_0/h = \num{3.04e-19}/\num{6.63e-34} \approx
\qty{4.6e14}{Hz}$; $\lambda_0 = c/\nu_0 \approx \qty{654}{nm}$. Pada
\qty{633}{nm}, $E = 1240/633 = \qty{1.96}{eV} > \qty{1.9}{eV}$: ya,
nyaris saja, dengan $E_{k,\max} \approx \qty{0.06}{eV} \approx \qty{1e-20}{J}$.
\end{solution}

\begin{solution}{exo:g12:quantum-world:4}
$E_3 - E_2 = -1.51 + 3.40 = \qty{1.89}{eV}$;
$\lambda = 1240/1.89 \approx \qty{656}{nm}$: garis merahnya.
\end{solution}

\begin{solution}{exo:g12:quantum-world:5}
$\nu_0 \approx \qty{5.6e14}{Hz}$; $W_0 = h\nu_0 \approx
\qty{3.7e-19}{J} \approx \qty{2.3}{eV}$ (perpotongan $-W_0$-nya).
Kemiringannya adalah \omterm{def:g12:quantum-world:photon}{tetapan Planck} $h$.
\end{solution}

\begin{solution}{exo:g12:quantum-world:6}
$E = 1240/400 = \qty{3.10}{eV}$; $E_{k,\max} = 3.10 - 2.3 =
\qty{0.80}{eV} = \qty{1.3e-19}{J}$; $U_s = E_{k,\max}/e =
\qty{0.80}{V}$.
\end{solution}

\begin{solution}{exo:g12:quantum-world:7}
Ketika intensitasnya digandakan: arusnya berlipat dua (\omterm{def:g12:quantum-world:photon}{fotonnya} dua kali lipat), $U_s$
tidak berubah (\omterm{def:g10:energy-conservation:energy}{energi} tiap \omterm{def:g12:quantum-world:photon}{fotonnya} sama). Ketika \omterm{def:g12:oscillators-and-time:oscillator}{frekuensinya} dinaikkan: $U_s$ naik,
$E_{k,\max} = h\nu - W_0$. Gambaran \omterm{def:g10:signals-and-waves:wave}{gelombangnya} mengikat \omterm{def:g10:energy-conservation:energy}{energi} elektronnya
pada intensitasnya, bukan pada \omterm{def:g12:oscillators-and-time:oscillator}{frekuensinya} --- persis yang \emph{tidak}
teramati.
\end{solution}

\begin{solution}{exo:g12:quantum-world:8}
Dari $n = 1$: \qty{13.6}{eV}, $\lambda = 1240/13.6 \approx
\qty{91}{nm}$, ultraungu jauh. Dari $n = 2$: \qty{3.40}{eV},
$\lambda \approx \qty{365}{nm}$, ultraungu dekat.
\end{solution}

\begin{solution}{exo:g12:quantum-world:9}
Dari \omterm{def:g12:quantum-world:levels}{keadaan dasarnya} loncatan terkecil berongkos $E_2 - E_1 =
\qty{10.2}{eV}$ (\qty{122}{nm}): setiap garis serapannya berada di
ultraungu, tidak ada yang tampak. Di dalam atmosfer bintang yang panas tumbukannya menjaga
sebagian \omterm{def:g11:fundamental-interactions:ladder}{atomnya} di $n = 2$, dan dari sana \omterm{def:g12:quantum-world:photon}{foton} Balmer yang tampak
(\qtyrange{1.89}{3.40}{eV}) dapat terserap.
\end{solution}

\begin{solution}{exo:g12:quantum-world:10}
(a) $E_k = \qty{1.6e-17}{J}$, $p = \sqrt{2m_e E_k} =
\qty{5.4e-24}{kg.m/s}$, $\lambda = h/p \approx \qty{1.2e-10}{m}$ ---
sebesar \omterm{def:g11:fundamental-interactions:ladder}{atom}: hablurnya melenturkannya. (b) $p = \num{1.0e-4} \times 1.0 =
\qty{1.0e-4}{kg.m/s}$, $\lambda \approx \qty{6.6e-30}{m}$ --- $10^{20}$
kali lebih kecil daripada sebuah \omterm{def:g11:fundamental-interactions:ladder}{atom}: \omterm{def:g10:signals-and-waves:wave}{gelombang} lalatnya selamanya tak terlihat.
\end{solution}

\begin{solution}{exo:g12:quantum-world:11}
Satu \omterm{def:g12:quantum-world:photon}{foton}: $\num{1.99e-25}/\num{5.05e-7} = \qty{3.94e-19}{J}$; sembilan puluh:
$\approx \qty{3.5e-17}{J}$. Nyamuknya: $\tfrac12 \times \num{2.5e-6}
\times 0.50^2 = \qty{3.1e-7}{J}$ --- $10^{10}$ kali lebih besar. Matanya
bekerja dalam seratus butir dari batas mutlaknya: hampir sebuah pencacah
\omterm{def:g12:quantum-world:photon}{foton}.
\end{solution}

\begin{solution}{exo:g12:quantum-world:12}
$\nu_1 = \qty{7.41e14}{Hz}$, $\nu_2 = \qty{5.49e14}{Hz}$;
$h = e\,\Delta U_s/\Delta\nu = \num{1.60e-19} \times 0.79 /
\num{1.91e14} \approx \qty{6.6e-34}{J.s}$. $W_0 = h\nu_2 - eU_{s,2} =
\num{3.63e-19} - \num{0.64e-19} = \qty{3.0e-19}{J} \approx
\qty{1.9}{eV}$: sesium.
\end{solution}

\begin{solution}{exo:g12:quantum-world:13}
$1240/486 = \qty{2.55}{eV} = E_4 - E_2 = 3.40 - 0.85$: transisi
$4 \to 2$. $1240/434 = \qty{2.86}{eV} = E_5 - E_2 = 3.40 - 0.54$:
transisi $5 \to 2$.
\end{solution}

\begin{solution}{exo:g12:quantum-world:14}
Tampang lintangnya $\pi r^2 \approx \qty{3.1e-20}{m^2}$, \omterm{def:g11:work-of-force:power}{daya} yang terpungut
$\qty{3.1e-26}{W}$; $t = \num{3.7e-19}/\num{3.1e-26} \approx
\qty{1.2e7}{s}$ --- empat bulan, lawan di bawah senanosekon yang teramati.
\omterm{def:g10:energy-conservation:energy}{Energinya} tidak dapat tersebar pada \omterm{def:g12:mechanical-waves:wavefront}{muka gelombangnya}: ia berjalan
terpusat, dalam butiran.
\end{solution}

\begin{solution}{exo:g12:quantum-world:15}
$E_k = \qty{8.0e-16}{J}$, $p = \sqrt{2m_e E_k} =
\qty{3.8e-23}{kg.m/s}$, $\lambda = h/p \approx \qty{1.7e-11}{m} =
\qty{17}{pm}$ --- sekitar $\num{3e4}$ kali lebih pendek daripada \qty{550}{nm}.
Pelenturan mengaburkan rincian yang lebih kecil daripada $\sim\lambda$, sehingga menyusutkan
$\lambda$ sebesar $10^4$ membeli $10^4$ pada \omterm{def:g11:work-of-force:power}{daya} pilahnya.
\end{solution}

\begin{solution}{pb:g12:quantum-world:1}
\textbf{1.} $E = hc/\lambda = \num{1.99e-25}/\num{5.50e-7} =
\qty{3.62e-19}{J} = \qty{2.26}{eV}$.

\textbf{2.} $N = \num{1.0e3}/\num{3.62e-19} \approx \num{2.8e21}$
\omterm{def:g12:quantum-world:photon}{foton} tiap sekon.

\textbf{3.} $P = 0.20 \times \qty{1.0e3}{W} = \qty{200}{W}$;
$E = 200 \times 5.0 = \qty{1.0}{kW.h}$ ($\qty{3.6e6}{J}$).

\textbf{4.} Dengan $10^{21}$ butir tiap sekon, kedip statistiknya
luar biasa kecil: kebutirannya merata-rata menjadi arus yang mulus
sempurna, seperti pasir yang dituang cepat mengalir bagaikan air.

\textbf{5.} Kerusakannya menuntut satu butir \qtyrange{3}{4}{eV}: butir ultraungu
memenuhinya, butir \omterm{def:g10:light-spectra:wavelength}{inframerah} ($\sim \qty{0.1}{eV}$) tidak pernah ---
banyaknya tidak dapat menggantikan pukulan tiap butir.

\textbf{6.} $U_s$ adalah \omterm{def:g11:circuits-and-power:voltage}{tegangan} balik yang tepat menghentikan elektron
yang tercepat: $eU_s = E_{k,\max} = h\nu - W_0$.

\textbf{7.} $\nu = c/\lambda$: $8.22$, $7.41$, $6.88$ dan $5.49$
$\times\,10^{14}\,\unit{Hz}$.

\textbf{8.} $eU_s = h\nu - W_0$ bersifat afin pada $\nu$ dengan kemiringan $h$.
Kemiringan berturut-turutnya: $(0.94 - 0.37)/1.39 = 0.41$ dan
$(1.50 - 0.94)/1.34 = 0.42$ \unit{eV} tiap $10^{14}\,\unit{Hz}$ ---
ketiga titiknya segaris.

\textbf{9.} $h = e(U_{s,1}-U_{s,4})/(\nu_1-\nu_4) = \num{1.60e-19}
\times 1.13/\num{2.72e14} = \qty{6.64e-34}{J.s}$.

\textbf{10.} Pada \qty{546}{nm}: $W_0 = h\nu - eU_s = \num{3.64e-19} -
\num{0.59e-19} = \qty{3.05e-19}{J} = \qty{1.90}{eV}$;
$\nu_0 = W_0/h = \qty{4.60e14}{Hz}$; $\lambda_0 \approx \qty{653}{nm}$.

\textbf{11.} Hanya $\lambda < \qty{653}{nm}$ yang melontarkan: daerah tampak dari
ungu sampai merah, dan ultraungunya. Seluruh sisi \omterm{def:g10:light-spectra:wavelength}{inframerahnya} --- kira-kira separuh
\omterm{def:g10:energy-conservation:energy}{energi} Matahari --- tak berguna bagi katoda ini.

\textbf{12.} $(6.64 - 6.63)/6.63 \approx 0.2\%$ --- pengukuran dua
permil sebuah tetapan semesta di laboratorium fisika sekolah.

\textbf{13.} $E_2 = \qty{-3.40}{eV}$, $E_3 = \qty{-1.51}{eV}$,
$E_4 = \qty{-0.85}{eV}$, $E_5 = \qty{-0.54}{eV}$.

\textbf{14.} $3 \to 2$: \qty{1.89}{eV}; $4 \to 2$: \qty{2.55}{eV};
$5 \to 2$: \qty{2.86}{eV}.

\textbf{15.} $\lambda = 1240/E$: \qty{656}{nm} (merah), \qty{486}{nm}
(biru kehijauan), \qty{434}{nm} (ungu).

\textbf{16.} \omterm{def:g10:energy-conservation:energy}{Energi} \omterm{def:g11:fundamental-interactions:ladder}{atomnya} \omterm{def:g12:quantum-world:levels}{terkuantum}, sehingga ia hanya dapat memancarkan
selisih yang diskret $E_p - E_n$: sebuah tangga membuat garis, bukan
pelangi.

\textbf{17.} Tidak: \qty{500}{nm} berarti \qty{2.48}{eV}, yang tidak cocok dengan
selisih tingkat mana pun (dari $E_2$ ke atas memberikan $1.89$, $2.55$, $2.86$,
\qty{3.40}{eV}; dari $E_1$ segalanya berongkos sekurang-kurangnya \qty{10.2}{eV})
--- gasnya bening di sana.

\textbf{18.} $eU = \tfrac12 m_e v^2$: $v = \sqrt{2 \times
\num{1.60e-19} \times \num{1.0e4}/\num{9.11e-31}} \approx
\qty{5.9e7}{m/s} \approx 0.20\,c$ --- perlakuan klasiknya dapat diterima.

\textbf{19.} $p = m_e v = \qty{5.4e-23}{kg.m/s}$;
$\lambda = h/p = \qty{1.2e-11}{m} = \qty{12}{pm}$.

\textbf{20.} Perolehannya $= \num{5.50e-7}/\num{1.23e-11} \approx \num{4.5e4}$.
Kedua bilangan akhir pekannya: $h = \qty{6.64e-34}{J.s}$, yang terukur sampai
$0.2\%$ dengan beberapa tapis dan sebuah voltmeter, dan sebuah mikroskop yang
empat puluh lima ribu kali lebih tajam karena elektron bergelombang.
\end{solution}
